\exercisesection{2}

\begin{enumerate}
\def\labelenumi{\arabic{enumi}.}
\tightlist
\item
  \begin{enumerate}
  \def\labelenumii{(\alph{enumii})}
  \tightlist
  \item
    证明若 \(\mathfrak{a}\) 是一个 \(\neq 0\) 的 Dedekind 整环 \(\mathbf{A}\) 中的理想，则环 \(\mathbf{A} / \mathfrak{a}\) 是主理想环（代数第七章 § 1 习题 5）（注意 \(\mathbf{A} / \mathfrak{a}\) 是半局部的，并仿照第 2 节命题 1 论证）。再次推出 \(\mathbf{A}\) 的每个分式理想至多由两个元素生成（参见习题 11 (a)）。
  \end{enumerate}
\end{enumerate}

\begin{enumerate}
\def\labelenumi{(\alph{enumi})}
\setcounter{enumi}{1}
\tightlist
\item
  在多项式环 \(\mathbf{A} = k[\mathbf{X},\mathbf{Y}]\) 中（其基域为 \(k\)），令 \(\mathfrak{m}\) 为理想 \(\mathbf{AX} + \mathbf{AY}\)。证明对每个整数 \(n\)，理想 \(\mathfrak{m}^n\) 的最小生成元数为 \(n + 1\)。
\end{enumerate}

\begin{enumerate}
\def\labelenumi{\arabic{enumi}.}
\setcounter{enumi}{1}
\tightlist
\item
  \begin{enumerate}
  \def\labelenumii{(\alph{enumii})}
  \tightlist
  \item
    设 \(\mathfrak{a}\)、\(\mathfrak{b}\) 为 Dedekind 整环 \(\mathbf{A}\) 的两个整理想；证明存在一个整理想 \(\mathfrak{c}\)，它与 \(\mathfrak{b}\) 互素，并使得 \(\mathfrak{ac}\) 为主理想（使用第一节第 5 节命题 9）。
  \end{enumerate}
\end{enumerate}

\begin{enumerate}
\def\labelenumi{(\alph{enumi})}
\setcounter{enumi}{1}
\tightlist
\item
  设 \(\mathfrak{a}, \mathfrak{b}\) 为 \(\mathbf{A}\) 的两个整理想；证明存在 \(x \neq 0\) 属于 \(\mathbf{A}\) 的分式域，使 \(x\mathfrak{a}\) 是 \(\mathbf{A}\) 的一个与 \(\mathfrak{b}\) 互素的整理想（同法）。推出模 \(\mathfrak{a}/\mathfrak{ab}\) 同构于 \(\mathbf{A}/\mathfrak{b}\)。
\end{enumerate}

\begin{enumerate}
\def\labelenumi{\arabic{enumi}.}
\setcounter{enumi}{2}
\item
  设 A 为 Dedekind 整环，K 为其分式域，P 为 A 的素理想 \(\neq 0\) 的集合；对每个 \(p \in P\)，令 \(v_{p}\) 为 K 上相应的本性赋值。对 K 的每个子 A-模 M 以及每个 \(p \in P\)，记 \(v_{p}(M) = \inf v_{p}(x)\)（取于 \(\overline{\mathbf{R}}\) 中）；若 \(M \neq 0\)，则 \(v_{p}(M) < +\infty\) 对所有 \(p \in P\) 成立，且 \(v_{p}(M) \leqslant 0\) 对几乎所有 \(p \in P\) 成立。若 M、N 是 K 的两个子 A-模，证明关系 \(M \subset N\) 等价于 \(v_{p}(N) \leqslant v_{p}(M)\) 对所有 \(p \in P\) 成立（使用第一节第 5 节命题 9）。反之，对于每族 \((v_{p})_{p \in P}\) 的元素，每个元素等于某个整数或 \(-\infty\)，且 \(v_{p} \leqslant 0\) 对几乎所有 \(p \in P\) 成立，证明存在唯一的 K 的子 A-模 M，使 \(v_{p}(M) = v_{p}\) 对所有 \(p \in P\) 成立。
\item
  设 A 为 Dedekind 整环，K 为其分式域。若 L 是 K 的子域，且 A 对 \(A \cap L\) 为整，则 \(A \cap L\) 是 Dedekind 整环（先证明 \(A \cap L\) 是 Krull 整环，然后使用第五章 §2 第 1 节定理 1 证明 \(A \cap L\) 中每个素理想都是极大理想）。
\item
  \begin{enumerate}
  \def\labelenumii{(\alph{enumii})}
  \tightlist
  \item
    设 k 为域，K 为有理函数域 \(k(\mathbf{X}, \mathbf{Y})\)，即 k 上两个不定元的有理函数域；A 为一个不定元上的多项式环 K{[}Z{]}，它是主理想整环。令 L 为 \(k(\mathbf{Z}, \mathbf{X} + \mathbf{Y}\mathbf{Z})\) 的子域 \(k(\mathbf{X}, \mathbf{Y}, \mathbf{Z})\)；证明 \(A \cap L\) 不是 Dedekind 整环（证明这个环中存在非零且不是极大理想的素理想 \(\neq 0\)）。
  \end{enumerate}
\end{enumerate}

\begin{enumerate}
\def\labelenumi{(\alph{enumi})}
\setcounter{enumi}{1}
\tightlist
\item
  设 A 为非主理想整环的 Dedekind 整环，且其理想类群 C(A) 是有限的 *（有限代数扩张 Q 的整数环具有此性质）。* 令 \((\mathfrak{a}_{j})_{1 \leqslant j \leqslant r}\) 为 A 的分式理想群 I(A) 模去 \(\neq 0\) 的分式主理想子群的代表系，该代表系由 A 的整理想组成；令 \(p_{i} (1 \leqslant i \leqslant s)\) 为至少整除某个 \(a_{j}\) 的素理想。令 S（分别为 T）
\end{enumerate}

表示由 \(x \in \mathbf{A}\) 满足 \(v_{\mathfrak{p}_i}(x) = 0\) 对所有 \(i\) 组成的乘性集（分别地，由 1 以及满足 \(x \in \mathbf{A}\) 且 \(v_{\mathfrak{q}}(x) = 0\) 的、其中 \(\mathfrak{q}\) 异于 \(\mathfrak{p}_i\) 的所有素理想，并且满足 \(v_{\mathfrak{p}_i}(x) \geqslant 1\) 对所有 \(i\) 的元素组成；假设蕴含 \(\mathbf{T}\) 不退化为 1）。证明 \(S^{-1}\mathbf{A}\) 和 \(T^{-1}\mathbf{A}\) 是主理想整环，且 \(A = (S^{-1}A) \cap (T^{-1}A)\)。

\begin{enumerate}
\def\labelenumi{\arabic{enumi}.}
\setcounter{enumi}{5}
\item
  证明在 Dedekind 整环中，准素理想、不可约理想、素理想（第四章 § 2 习题 33）、拟素理想（第四章 § 2 习题 34）以及素理想的幂这几个概念彼此相同。
\item
  设 A 为 Noether 整环。证明下列性质彼此等价：
\end{enumerate}

\((\alpha)\) A 是 Dedekind 整环。

(β) 对 A 的每个极大理想 m，不存在异于 m 和 \(m^{2}\) 的理想 a，使 \(m^{2} \subset a \subset m\)。

\((\gamma)\) 对每个极大理想 \(\mathfrak{m}\)（属于 \(\mathbf{A}\)），关于 \(\mathfrak{m}\) 的准素理想按包含关系全序。

(δ) 对 A 的每个极大理想 m，关于 m 的每个准素理想都是素理想的乘积。

（证明性质 \((\gamma)\) 和 \((\delta)\) 各自蕴含 \((\beta)\)；然后证明 \((\beta)\) 蕴含 \(A_{m}\) 是域或离散赋值环（参见第六章 § 3 第 5 节命题 9）。）

举出一个满足性质 \((\beta)\)、\((\gamma)\) 和 \((\delta)\) 但不是赋值环的局部整环的例子（参见第六章 § 3 习题 7）。

¶ 8. 设 A 为整环，其中每个素理想 ≠0 都是可逆的。证明 A 是 Dedekind 整环。（先证明 A 的每个素理想 p' ≠ 0 都是极大理想，注意若 p 是满足 p ≠ p' 且 p ⊃ p' 的素理想，则 p': p = p'（第二章 § 1 习题 8 (b)），另一方面 p': p = p'p⁻¹，从而得到矛盾。推出 A 是 Noether 环（第二章 § 1 习题 6 (b)），最后使用第二章 § 5 第 6 节定理 4，证明对 A 的每个极大理想 m，A\_m 是域或离散赋值环。）

\begin{enumerate}
\def\labelenumi{\arabic{enumi}.}
\setcounter{enumi}{8}
\tightlist
\item
  设 A 为整环。
\end{enumerate}

\begin{enumerate}
\def\labelenumi{(\alph{enumi})}
\item
  设 \((\mathfrak{p}_i)_{1\leqslant i\leqslant m},(\mathfrak{p}_j')_{1\leqslant j\leqslant n}\) 为两个有限的可逆素理想族，且 \(\mathfrak{p}_1\mathfrak{p}_2\dots \mathfrak{p}_m = \mathfrak{p}_1'\mathfrak{p}_2'\dots \mathfrak{p}_n'\)。证明 \(m = n\)，并证明存在一个置换 \(\pi\)，它是 \([1,n]\) 上的置换，使得 \(\mathfrak{p}_i' = \mathfrak{p}_{\pi (i)}\) 对所有 \(i\)（满足 \(1\leqslant i\leqslant n\)）成立。
\item
  设 \(\mathfrak{p}\) 为 \(\mathbf{A}\) 的素理想，且 \(a\) 属于 \(\mathbf{A} - \mathfrak{p}\)；若 \(\mathfrak{p} \subset \mathbf{A}a + \mathfrak{p}^2\)，证明 \(\mathfrak{p} = \mathfrak{p}(\mathbf{A}a + \mathfrak{p})\)；推出若 \(\mathfrak{p}\) 可逆，则 \(\mathbf{A}a + \mathfrak{p} = \mathbf{A}\)。
\item
  证明若 A 的每个理想都是 A 的素理想之积（事先不要求唯一），则 A 是 Dedekind 整环。（先证明每个可逆素理想 p 都是极大理想：考虑 \(a \in A - p\)，将 \(Aa + p\) 和 \(Aa^{2} + p\) 分解为素理想之积，并应用 (a)
\end{enumerate}

在环 A/p 中证明 \(Aa^{2} + p = (Aa + p)^{2}\)；然后使用 (b)。再证明每个 \(p \neq 0\) 的素理想都是可逆的：将含有 \(b \in p\) 的主理想 Ab 分解为素理想之积，注意该积的因子都是可逆的，并应用上述结论证明 p 必等于这些因子中的一个。借助习题 8 得出结论。）

¶ 10. 设 A 为一个环，其中每个理想都是素理想之积。

\begin{enumerate}
\def\labelenumi{(\alph{enumi})}
\item
  证明对 A 的每个素理想 p，A/p 是 Dedekind 整环（习题 9）。
\item
  证明 A 的极小素理想  的集\(\mathfrak{p}_i\)\(\mathbf{A}\)证明 A 的极小素理想  的集合是有限的（将 (0) 分解为素理想之积）。
\item
  假设 \(\mathbf{A} / \mathfrak{p}_i\) 不是域。证明若 \(y \in \mathfrak{p}_i\)，则对所有 \(x \in \mathbf{A} - \mathfrak{p}_i\)，存在 \(z \in \mathfrak{p}_i\)，使 \(y = zx\)。（在 A 中考虑 Ax 和 Ax + Ay 分解为素理想之积，并考虑它们在 \(\mathbf{A} / \mathfrak{p}_i\) 中的相应分解。）推出对所有非零 \(y \in \mathfrak{p}_i\)，\(Ay = \mathfrak{p}_i\)（考虑 \(\mathfrak{p}_i / Ay\)，将 Ay 分解为素理想之积）；最后证明 \(\mathfrak{p}_i = \mathfrak{p}_i^2\)，从而在 \(\mathfrak{p}_i\) 中存在幂等元 \(e_i\)，使 \(\mathfrak{p}_i = Ae_i\)（参见第二章 § 4 习题 15）。于是 A 是环 \(Ae_i\) 与环 \(\mathbf{A}(1 - e_i)\) 的直积，后者同构于 Dedekind 整环 \(\mathbf{A} / \mathfrak{p}_i\)。
\item
  现在假设所有 \(p_{i}\) 都是极大理想，于是 A 是形如 \(A/p_{i}^{v_{i}}\) 的环的直积（第二章 § 1 第 2 节命题 5），因此只需考察 A 为准素环或 A 的唯一素理想 p 为幂零的情形。证明在此情形下，该假设蕴含 A 是主理想环（代数第七章 § 1 习题 5、6）。
\item
  由 (c) 和 (d) 推出 A 是有限个 Dedekind 整环与一个主理想环的乘积。
\end{enumerate}

\begin{enumerate}
\def\labelenumi{\arabic{enumi}.}
\setcounter{enumi}{10}
\tightlist
\item
  设 K 为域\((v_{\iota})_{\iota\in I}\)， 为 K 上满足 §1 \((\mathbf{AK}_{\mathbf{I}})\)\((\mathbf{AK}_{\mathbf{III}})\)设 K 为域， 为 K 上满足 §1 第\((n_{h})_{1\leqslant h\leqslant r}\)\(\geqslant0\) \((\iota_{h})_{1\leqslant h\leqslant r}\)为域， 为 K 上满\((a_{h})_{1\leqslant h\leqslant r}\)\(x\in K\)\(v_{\iota_{h}}(x-a_{h})\geqslant n_{h}\)\(1\leqslant h\leqslant r\)\(v_{\iota}(x)=0\)设 K 为域，\(\iota_{h}\) 为 K 上满足 §1 第\(v_{\iota}\)\((v_{\iota})_{\iota\in I}\) 3 节的  和  条件的离散赋值族，并且对每个整数 r、每族  的非负整数、每族\(\neq0\)  的 I 中彼此不同的元素以及每族  的 K 中元素，都存在 ，使得 （），并且对不同于各  的 ι 有 。证明赋值环  的交 A 是 Dedekind 整环，其分式域为 K，且  是 A 的本性赋值族。（使用第一节习题 7；然后证明两个不同的高度为 1 的素理想互素，并推出 A 的每个素理想  都是极大理想。）
\end{enumerate}

¶ 12. 设 A 为整环，K 为其分式域。证明下列性质彼此等价：

\((\alpha)\) 对 A 的每个素理想 \(\mathfrak{p}\)，局部环 \(A_{\mathfrak{p}}\) 都是赋值环。

(β) 对 A 的每个极大理想 m，局部环 \(A_{m}\) 都是赋值环。

 A 中每个有限生成理\(\neq0\)想  都是可逆的（因而有限\(\neq0\)生成分式理想  构成一个群）。

(δ) 每个无挠有限生成 A-模都是投射模。

(ε) 对所有 \(x \in \mathbf{K}\)，分式理想 \(\mathbf{A} + \mathbf{A}x\) 都是可逆的。

(ζ) 对所有 \(x \neq 0\) 的 \(\mathbf{K}\) 中元素，存在 \(y \in \mathbf{K}\)，使 \(y \equiv 0 \pmod{1}\)、\(y \equiv 0 \pmod{x}\) 且 \(y \equiv 1 \pmod{1 - x}\)。

\((\eta)\) 若 A 的理想有限族为 \((a_{i})\)，则同余方程组 \(x \equiv c_{i} \pmod{a_{i}}\) 有解的充要条件是对每一对指标 i、j 都有 \(c_{i} \equiv c_{j} \pmod{(a_{i} + a_{j})}\)（“中国剩余定理”）。

(θ) 若 \(a, b, c\) 是三个 \(A\) 的理想，则 \(a \cap (b + c) = a \cap b + a \cap c\)。

\begin{enumerate}
\def\labelenumi{(\roman{enumi})}
\item[(ι)]
(ι) 若 \(a, b, c\) 是三个 \(A\) 的理想，则 \(a + (b \cap c) = (a + b) \cap (a + c)\)。
\item[(κ)]
(κ) A 是整闭的，且 A 的每个有限生成理想都是除性的。
\end{enumerate}

(λ) A 是整闭的，并且对属于 \(z \neq 0\)\(\mathbf{K}\)) A \(x, y\)\(\mathbf{A}\)) A\(z = x + yz^2\) 是整闭的，并且对属于  的 ，存在属于  的 ，使 。

(μ) A 是整闭的，并且若 a、b、c 是三个有限生成分式理想 \(\neq0\)，关系 ab = ac 蕴含 b = c。

\begin{enumerate}
\def\labelenumi{(\alph{enumi})}
\setcounter{enumi}{21}
\tightlist
\item
  对每个有限生成 A-模 M，M 的挠子模是 M 的直和因子。
\end{enumerate}

(σ) 每个满足 A ⊂ B ⊂ K 的环 B 都是整闭整环。

于是称 A 为 Prüfer 整环（或 Prüferian 整环）。

（证明 \((\alpha)\)、\((\beta)\)、\((\gamma)\) 和 \((\delta)\) 的等价性，使用第二章 § 5 第 2 节定理 1 和第 6 节定理 4。为证明 \((\varepsilon)\) 蕴含 \((\gamma)\)，使用整环中三个分式理想之间的恒等式 \((a + b)(b + c)(c + a) = (a + b + c)(ab + bc + ca)\)。为证明 \((\zeta)\) 蕴含 \((\eta)\)，对 \(a_i\) 的个数作归纳。\((\eta)\)、\((\theta)\) 和 \((\iota)\) 的等价性已在代数第六章 § 1 习题 25 中证明。证明 \((\lambda)\) 蕴含 \((\varepsilon)\)：注意 \((\lambda)\) 蕴含关系

\[
x (\mathbf {A} + \mathbf {A} z) \subset z (\mathbf {A} + \mathbf {A} z),
\]

使用§ 1 的习题 6 (b)，并注意

\[
(\mathbf {A} + \mathbf {A} z) (\mathbf {A} y + \mathbf {A} x / z) = \mathbf {A}.
\]

证明 \((\mu)\) 蕴含 \((\lambda)\)，注意恒有

\[
z (\mathrm{A} + \mathrm{Az}) \subset (\mathrm{A} + \mathrm{Az} ^ {2}) (\mathrm{A} + \mathrm{Az}).
\]

证明 \((\kappa)\)\((\lambda)\) 蕴含 ，注意每个分式\(A t\)\(A\)\(A z^2\)明  \(A z\)蕴\((\nu)\)\((\gamma)\) 蕴含 ，注意每个分式主理\(\mathfrak{a} \neq 0\)\(A\)明  蕴含 \(c \in \mathfrak{a}(A: \mathfrak{a})\)，注意每个\(\mathfrak{b} = c \mathfrak{a}\)\((\sigma)\)\((\beta)\) 蕴含 ，注\(A\)意每\(z \in K - A\)\(z^{-1} \in A\)蕴含 ，注意每个分\(z \in A[z^2]\)式主理想  都包含  和 ，也包含 。为证明  蕴含 ，对非零有限生成整理想  的 ，取非零元素 ，并对理想  应用第一节习题 32。为证明  蕴含 ，将问题化为  是局部环的情形，取 ，证明 ，注意该假设蕴含 。

¶ 13. (a) 在 Prüfer 整环 A 中，设 \(a\)、\(b\) 是两个有限生成理想；证明：若 \(b \notin a\)，则存在有限生成理想 \(c\)，使得 \(a \subset c\)、\(a \neq c\) 且 \(bc \subset a\)（考虑理想 \(a + b\)）。

\begin{enumerate}
\def\labelenumi{(\alph{enumi})}
\setcounter{enumi}{1}
\item
  由 (a) 推出：若 a 是 Prüfer 整环 A 中的有限生成理想，则在环 A/a 中，每个不是零因子的元素都是可逆的。特别地，如果素理想是极大理想，则 A 的素理想不可能是有限生成的。
\item
  证明在 Prüfer 整环 A 中，两个互不包含于对方的素理想 p、p' 互素（对 A 的每个极大理想 m，考虑理想 \(pA_{m} + p'A_{m}\)）。
\item
  在 Prüfer 整环 A 中，设 \(\mathfrak{a}\) 是一个有限生成理想。要使 \(\mathfrak{a}\) 成为本原理想（第四章，§ 2，习题 33），必要且充分条件是 \(\mathfrak{a}\) 只包含于 A 的一个极大理想中（使用 (b)）；此时 \(\mathfrak{a}\) 是不可约且拟素的（第四章，§ 2，习题 34），所以对于有限生成理想，这三个概念彼此等价。要使 \(\mathfrak{a}\) 成为准素理想，必要且充分条件是它只包含于一个素理想（必为极大理想）\(\mathfrak{m}\) 中；要使 \(\mathfrak{a}\) 成为强准素理想（第四章，§ 2，习题 27），必要且充分条件还包括 \(\mathfrak{mA}_{\mathfrak{m}}\) 是主理想。
\end{enumerate}

\begin{enumerate}
\def\labelenumi{\arabic{enumi}.}
\setcounter{enumi}{13}
\item
  设 A 是 Prüfer 整环。证明：对于 A-模 M，平坦的必要且充分条件是无挠（使用习题 12 (δ)）。推出 A 是凝聚环（第一章，§ 2，习题 12）。
\item
  \begin{enumerate}
  \def\labelenumii{(\alph{enumii})}
  \tightlist
  \item
    若 A 是 Prüfer 整环，则对于 A 的每个素理想 p，A/p 也是 Prüfer 整环。
  \end{enumerate}
\end{enumerate}

\begin{enumerate}
\def\labelenumi{(\alph{enumi})}
\setcounter{enumi}{1}
\item
  若 A 是 Prüfer 整环，则对于 A 的每个不含 0 的乘法子集 S，\(S^{-1}A\) 也是 Prüfer 整环。
\item
  设 K 是一个域，\((\mathbf{A}_{\lambda})\) 是 K 的子环组成的一个非空右有向族。证明：若各个 \(A_{\lambda}\) 都是 Prüfer 整环，则它们的并 A 也是 Prüfer 整环（使用习题 12 (ζ)）。
\end{enumerate}

\begin{enumerate}
\def\labelenumi{\arabic{enumi}.}
\setcounter{enumi}{15}
\tightlist
\item
  设 A 是 Prüfer 整环，K 是它的分式域，L 是 K 的一个（有限或非有限）代数扩张；证明 A 在 L 中的整闭包 A' 是 Prüfer 整环。（取 \(x \in L\)，设 \(a_{0}X^{n} + a_{1}X^{n-1} + \cdots + a_{n}\) 是 A{[}X{]} 中以 x 为根的多项式；若我们写成
\end{enumerate}

\[
a _ {0} \mathrm{X} ^ {n} + a _ {1} \mathrm{X} ^ {n - 1} + \dots + a _ {n} = (\mathrm{X} - x) (b _ {0} \mathrm{X} ^ {n - 1} + \dots + b _ {n - 1}),
\]

  证明 \(b_{i}\) 和 \(b_{i}x\) 属于 \(A'\)，使用 §1，习题 12 (a)。若 \(a = \sum_{i=0}^{n} A a_{i}\)，\(b = \sum_{j=0}^{n-1} B b_{j}\)，证明理想 \(b.B a^{-1}\) 是 \(B + B x\) 的逆理想。）

¶ 17. 设 A 是一个整环。

\begin{enumerate}
\def\labelenumi{(\alph{enumi})}
\item
  对于 A 成为 Bézout 整环（§ 1，习题 20），必要且充分条件是它是 Prüfer 整环且是伪 Bézout 整环（§ 1，习题 21）。
\item
  对于 A 成为 Dedekind 整环，必要且充分条件是它是 Prüfer 整环 Krull 整环（证明 A 的每个除性理想都是有限生成的，使用习题 12 (κ) 和 § 1，习题 11 (a)；推出 A 的每个素理想 \(\neq 0\) 都是极大理想（习题 13 (b)），然后推出 A 的每个素理想都是有限生成的，并借助第二章，§ 1，习题 6 得出结论）。
\item
  对于 A 成为 Dedekind 整环，必要且充分条件是它是 Prüfer 整环且强 Lasker（第四章，§ 2，习题 28）。（注意：若 A 是 Prüfer 整环且强 Lasker，则 A 的每个极大理想 m 都使得 \(A_{m}\) 是离散赋值环；进而对 A 中每个非零 \(x \in A\)，存在一个由 A 的有限个极大理想组成的乘积，包含于理想 Ax 中，由此得出 A 是 Krull 整环。）
\end{enumerate}

\begin{enumerate}
\def\labelenumi{\arabic{enumi}.}
\setcounter{enumi}{17}
\tightlist
\item
  设 A 是一个环，它是由一族右有向的子环 \(A_{\alpha}\) 的并，而这些子环都是 Dedekind 整环。假设：对于 \(A_{\alpha}\) 的每个素理想 p，存在 \(\beta \geqslant \alpha\)，使得 \(A_{\beta}\) 中至少有两个不同的素理想位于 p 之上。
\end{enumerate}

\begin{enumerate}
\def\labelenumi{(\alph{enumi})}
\item
  证明所有代数整数构成的环（\(\mathbf{Z}\) 在 \(\mathbf{C}\) 中的整闭包）满足上述条件（参见第五章，§ 2，习题 6）。
\item
  \(v\)设  是  上的真赋\(\mathbf{A}\)\(z \in \mathbf{A}\) 设 \(v(z) > 0\)\(\mathfrak{a}\)  是\(\mathbf{A}\) \(x\)\(v(x) \geqslant v(z)\) \(\mathbf{A}: \mathfrak{a} = \mathbf{A}\)是\(\mathfrak{a}\)  上的真赋值， 且 ，并\(m\)\(\mathbf{A}\)\(\mathfrak{a}\mathrm{A}_{m}\) 设  是  上的真赋值， 且\(\mathrm{A}_{\mathfrak{m}}\) ，并设  是  中由满足  的 x 组成的理想。证明 ，从而  不可逆，尽管对于  的每个极大理想 m， 是赋值环  中的主理想（参见第二章，§ 5，第 6 节，定理 4）。
\item
  由 (a) 和 (b) 得到一个完全整闭但不是 Krull 整环的 Prüfer 整环例子（参见习题 15 (c) 和 17 (b) 以及第五章，§ 1，习题 14）。
\end{enumerate}

¶ 19. 若整环 A 的 D(A)（§ 1，习题 11）的有限生成除子集合构成一个群，则称 A 为伪 Prüfer 整环。伪 Prüfer 整环是正则整闭的（§ 1，习题 29）。伪 Bézout 整环（§ 1，习题 21）是伪 Prüfer 整环；Prüfer 整环是伪 Prüfer 整环；Krull 整环是伪 Prüfer 整环。

\begin{enumerate}
\def\labelenumi{(\alph{enumi})}
\item
  通过例子证明，最后三个断言的逆命题不一定成立。
\item
  \(\theta\)设  是一个\(>0\)\(\Gamma\)设 \(\mathbf{R}^2\) 是一个  的无理数\((\alpha, \beta)\)\(\alpha \geqslant 0\)\(\beta \geqslant \theta \alpha\)设  是一个  的无理数， 是\(\Gamma\)群 ，其序由以\(\Gamma\)下方式给出：取满足 \(k\)\(\mathbf{Q}^2\)\(\mathbf{R}^2\) 设  是一个  的无理数， 是群 ，其序由以下方式给出：取满足  且  的有序对  作为正元素；有序群  是一个完备格。\(\mathrm{K}^*/\mathrm{U}\)\(\Gamma \cap \mathbf{Q}^2\) 的无理数，\(\Gamma\) 是群 ，其序由以下方式给出：取满足  且  的有序对  作为正元素；有序群  是一个完备格。设 B 是由 § 1，习题 22 的过程从  导出的伪主理想整环，设 A 是 B 与  的子群  在 k 上的代数的交。证明 A 是完全整闭整环，但不是伪 Prüfer 整环（注意：若 K 是 A 的分式域，U 是 A 的可逆元素群，则有序群  同构于 ，其序由  上的序诱导）。
\item
  若整环 A 是伪 Prüfer 整环，证明多项式整环 A{[}X{]} 是伪 Prüfer 整环（参见 § 1，习题 30 (c)）。
\item
  设 A 是伪 Prüfer 整环，K 是其分式域。证明 A 在 K 的一个代数扩张 L 中的整闭包 B 是伪 Prüfer 整环（采用习题 16 的方法，并使用 § 1，习题 29 (a)）。
\end{enumerate}

* (e) 由 § 1 的习题 30 (d) 推出一个例子：存在一个具有相同分式域的因子分解整环递增序列 \((\mathbf{A}_n)\)，其并不是正则整闭整环（因而更不可能是伪 Prüfer 整环）（在所引用的例子中，取 B 为离散赋值环）。 *

\begin{enumerate}
\def\labelenumi{\arabic{enumi}.}
\setcounter{enumi}{19}
\tightlist
\item
  设 A 是 Dedekind 整环，K 是其分式域，L 是 K 的有限代数扩张，B 是 A 在 L 中的整闭包，它是 Dedekind 整环。设 f 是 B 的一\(A \subset C \subset B\)\(p'\)\(B/p'\)\(A/(p' \cap A)\) Dedekind 整环，K 是其分式域，L 是\(p'\)\(p'\)设 A 是 Dedekind 整环，K 是其分式域，L 是 K 的有\(C_{0} = A + f\)\(A + f\)是 K 的有限代数扩张，B 是 A 在 L 中的整闭包，它是 Dedekind 整环。设 f 是 B 的一个理想。要存在一个环 C，使得 ，且 f 是 B 在 A 中的导子（第五章，§ 1，第 5 节），必要且充分条件是：对于 B 的每个包含 f 的素理想 ，若域  同构于 （剩余次数为 1 的素理想），则 f 与 p' 在 f 中的传递子 f:  在 A 中的交相等。（注意：如果存在这样的环 C，则 B 在环  中的导子也等于 f；因此 C 的存在等价于这样的事实： 不包含异于 f 且包含 f 的 B 理想。要证明后一条件等价于命题中的条件，将问题化为 A 是离散赋值环的情形。）
\end{enumerate}

¶ 21. (a) 设 A 是一个整环，f、g 是 A{[}X{]} 中的两个多项式，且 h = fg。记 a、b、c 分别为由 f、g、h 的系数生成的 A 的理想。证明：若 \(\deg(g) = n\)，则 \(a^{n+1}b = a^{n}c\)。（令 \(f(X) = \sum_{i=1}^{m} a_{i}X^{i}\)，\(g(X) = \sum_{j=1}^{n} b_{j}X^{j}\)；对于每个递增序列

\[
\sigma = (i _ {k}) _ {1 \leqslant k \leqslant n + 1}
\]

的整数 ，\(\leqslant m\)\(j\)，且对所有 ，令

\[
u _ {\sigma , j} = a _ {i _ {1}} a _ {i _ {2}} \dots a _ {i _ {j + 1}} b _ {j} a _ {i _ {j + 2}} \dots a _ {i _ {n + 1}};
\]

在集合  上取\(u_{\sigma,j}\)一个全序，\(j < j'\)\(u_{\sigma,j} < u_{\tau,j'}\)集合\(\sigma, \tau\)  上\(j\)\(u_{\sigma,j} \leqslant u_{\tau,j}\)集合  上\(\sigma \leqslant \tau\)取一个全序，使\([0, m]^{n+1}\)得当  时，对所有  都有 ；并且对每个 ，当且仅当在  上按字典序排列时 ，有 。然后对这个全序集作归纳论证。

\begin{enumerate}
\def\labelenumi{(\alph{enumi})}
\setcounter{enumi}{1}
\item
  由 (a) 推出：如果 A 是 Prüfer 整环（习题 12），则 c = ab。
\item
  取 A 为多项式环 \(\mathbf{Z}[\mathbf{Y}]\)\(f, g\)\(\mathrm{A}[\mathbf{X}]\)\(c \neq ab\) A 为多项式环 ，给出  中两个一次多项式  的例子，使得 。
\end{enumerate}

¶ 22. (a) 设 A 是一个 Noether 整环，其每个素理想  都是极大理想，因此每个理想 a  都能唯一写\(\prod_{i} q_{i}\)\(p_{i}\)er 整环，其每个素理想  都是极大理想，因此每个理想 a  都能唯一写成相对于包含 a 的不同素理想 p＿i 的准素理想之积 。

 \(\mathfrak{a}\) 设  是  中的\(\mathbf{A}\)\(\mathfrak{b}\)  \(\neq 0\)\(\mathbf{A}\) 设  是  中的可逆理\(\mathfrak{c}\)\(\mathfrak{b} + \mathfrak{c} = \mathbf{A}\)\(\mathfrak{ac}\)设  是  中的可逆\((\mathfrak{p}_i)_{1 \leqslant i \leqslant n}\)理想， 是 \(\mathbf{A}\)\(\mathfrak{r}_i = \prod_{j \neq i} \mathfrak{p}_j\)\(\mathfrak{a} = \sum_i \mathfrak{ar}_i\)\(\bigcap_i \mathfrak{ar}_i = \mathfrak{ap}_1\mathfrak{p}_2 \ldots \mathfrak{p}_n\)设  是  中的可逆理想， 是  中任意的  理想；证明\(x \in \mathfrak{a}\)\(x^{-1}\mathfrak{a} + \mathfrak{b} = \mathbf{A}\) 是\(\mathfrak{a}\)  中的可逆理想， 是  中任意的  理想；证明存在理想 ，使得 ，且  是主理想（注意：若  是  的有限个互异极大理想组成的族，且记 ，则 ，且 ；使用第二章，§ 1，第 2 节，命题 6，推出存在元素 ，使得 。）推出  由两个元素生成（参见习题 1）。

\begin{enumerate}
\def\labelenumi{(\alph{enumi})}
\setcounter{enumi}{1}
\tightlist
\item
  设 K 是一个域，A 是形式幂级数环 \(K[[T]]\) 的子环，由形如 \(a_{0} + T^{n}P(T)\) 的级数组成，其中 \(a_{0} \in K\)、\(P(T) \in K[[T]]\)，对给定整数 n 而言。证明 A 是一个 Noether 局部整环，只有一个素理想 \(m \neq 0\)，但 m 的生成系的最小基数是 n。
\end{enumerate}

